% Pandoc template for the fed_fraud documents.
% Based on proposal_gen/templates/detailed_proposal_template/detailed_1_classic.tex
% (Classic report): centered title page, table of contents, ruled section
% headings and the Geodesic house footer. Fonts are Latin Modern (lmodern).
% Two layouts: the full report (title page and contents) and, with
% -V brief, a compact one-page-brief header in the same style.
\documentclass[letterpaper,11pt]{article}
\usepackage[T1]{fontenc}
\usepackage[utf8]{inputenc}
\usepackage{lmodern}
\usepackage[letterpaper,top=0.9in,bottom=1.2in,left=1in,right=1in,footskip=0.6in]{geometry}
\usepackage[table]{xcolor}
\usepackage{geodesicbrand}
\usepackage{amsmath,amssymb}
\usepackage{enumitem}
\usepackage{array}
\usepackage{calc}
\usepackage{longtable}
\usepackage{makecell}
\usepackage[most]{tcolorbox}
\usepackage{titlesec}
\usepackage{needspace}
\usepackage{etoolbox}
\usepackage{hyperref}
\hypersetup{colorlinks=true,urlcolor=coolblack,linkcolor=coolblack,filecolor=geoblue,
  pdftitle={Justice},pdfauthor={\geofirm}}
\urlstyle{same}

\setlist{itemsep=2pt,leftmargin=1.4em}
\setlength{\parindent}{0pt}
\setlength{\parskip}{6pt}
\providecommand{\tightlist}{\setlength{\itemsep}{0pt}\setlength{\parskip}{0pt}}
\setcounter{secnumdepth}{0}

% Section style: house rule under each heading
\titleformat{\section}{\Large\bfseries\color{coolblack}}{\thesection}{0.6em}{}[{\color{coolblack}\titlerule[0.5pt]}]
\titleformat{\subsection}{\large\bfseries\color{geomid}}{\thesubsection}{0.6em}{}
\titleformat{\subsubsection}{\normalsize\bfseries\color{geomid}}{\thesubsubsection}{0.6em}{}
\titlespacing*{\section}{0pt}{10pt}{6pt}

% Tables are written by latex/filters.lua (navy header, banded rows)
\setlength{\LTpre}{6pt}\setlength{\LTpost}{6pt}

\begin{document}

% BRIEF HEADER
\thispagestyle{geodesic}
\begin{minipage}[c]{0.13\textwidth}\includegraphics[width=\linewidth]{icon.png}\end{minipage}\hfill
\begin{minipage}[c]{0.84\textwidth}
  {\textsc{\large \geofirm}}\par\vspace{2pt}
  {\small\color{geogray}Federal fraud expansion\enspace\textperiodcentered\enspace When: \compiletime}
\end{minipage}\par\vspace{6pt}
{\color{coolblack}\rule{\linewidth}{0.3mm}}\par\vspace{2pt}
{\LARGE\bfseries\color{coolblack} Justice}\par
\vspace{2pt}{\large\color{geomid} One-page brief}\par
\vspace{-2pt}{\color{coolblack}\rule{\linewidth}{0.3mm}}\par\vspace{4pt}
\pagestyle{geodesic}

\textbf{Priority:} 20 of 33 (score 2.9). \textbf{FY2026 outlays:} \$52B
(USAspending). \textbf{As of:} 2026-10-27. Public data only.

\hypertarget{bottom-line}{%
\section{Bottom line}\label{bottom-line}}

DOJ matters more as our customer than as a target: it runs the
whistleblower program behind america.gov/fraud, the Health Care Fraud
Unit, the Trade Fraud Task Force and the Procurement Collusion Strike
Force. Its own fraud risk is in grants (Office of Justice Programs, the
Crime Victims Fund passed through states) and in prison contracts. The
best DOJ product is a case-mining pipeline that turns its press releases
into labeled training data for every other brief.

\hypertarget{exposure}{%
\section{Exposure}\label{exposure}}

\begin{itemize}
\tightlist
\item
  \$52B in outlays, mostly salaries, prisons and grants to states and
  nonprofits.
\item
  Grant risk: subgrantees of victim and public safety grants are many
  and small, and DOJ OIG regularly finds unsupported costs.
\item
  No DOJ program is listed as high-risk for improper payments.
\end{itemize}

\hypertarget{public-data-we-can-use}{%
\section{Public data we can use}\label{public-data-we-can-use}}

\begin{itemize}
\tightlist
\item
  \textbf{DOJ press releases} (all U.S. Attorneys' offices), searchable
  and structured enough to extract defendants, schemes, programs and
  dollars.
\item
  \textbf{USAspending grants and subawards} for OJP, OVW and COPS.
\item
  \textbf{DOJ OIG audit reports} of grantees.
\item
  \textbf{FPDS} for Bureau of Prisons and other contracts.
\end{itemize}

\hypertarget{detection-plays}{%
\section{Detection plays}\label{detection-plays}}

\begin{enumerate}
\def\labelenumi{\arabic{enumi}.}
\tightlist
\item
  \textbf{Case mining (for all agencies):} extract defendant names,
  entities, addresses, schemes and agency from press releases; join to
  SAM.gov and USAspending to make labeled examples.
\item
  \textbf{Grant subrecipient risk:} subgrantees with prior OIG findings,
  shared officers across grantees, or award growth without performance
  data.
\item
  \textbf{Prison contract outliers:} commissary, medical and telecom
  vendors with prices far above peers.
\end{enumerate}

\hypertarget{first-prototype}{%
\section{First prototype}\label{first-prototype}}

A press-release case database: every fraud case since 2020 tagged by
agency and scheme and linked to award records. It becomes the label set
for every other prototype.

\hypertarget{limits}{%
\section{Limits}\label{limits}}

Press releases cover charged cases only, which biases labels toward
schemes DOJ already finds. Grant subaward reporting is incomplete.

\hypertarget{who-to-take-it-to}{%
\section{Who to take it to}\label{who-to-take-it-to}}

DOJ Criminal Division Fraud Section, Civil Division Commercial
Litigation (False Claims Act), Office of Justice Programs, and DOJ OIG.

\hypertarget{sources}{%
\section{Sources}\label{sources}}

\begin{itemize}
\tightlist
\item
  \href{https://www.justice.gov/fraud/whistleblower}{DOJ Whistleblower
  Rewards Program terms}
\item
  \href{https://www.justice.gov/news}{DOJ news releases}
\item
  \href{https://www.usaspending.gov/agency}{USAspending agency profiles
  (FY2026 outlays)}
\end{itemize}

\end{document}
